\documentclass[a4paper,12pt]{article}
\usepackage[utf8]{inputenc}
\usepackage{amsmath, amssymb}
\usepackage{graphicx}
\usepackage{booktabs}
\usepackage{geometry}
\geometry{margin=2.5cm}
\usepackage{auto-pst-pdf}
\usepackage{pstricks}
\usepackage{pst-labo}

\title{Forsøk: Loven om bevaring av masse}
\author{Analysefaget}
\date{\today}

\begin{document}

\maketitle

\section{Hensikt}
Å undersøke om den totale massen i et lukket system forblir konstant under fysiske og kjemiske prosesser. Dette demonstrerer \textbf{loven om bevaring av masse}.

\section{Teori}
Loven om bevaring av masse ble formulert av Antoine Lavoisier på slutten av 1700-tallet. Den sier at \emph{massen verken kan skapes eller ødelegges i en kjemisk reaksjon}. Atomene omorganiseres, men total mengde materie forblir den samme.

Ved kjernefysiske reaksjoner må man derimot ta hensyn til \emph{bevaring av masse og energi} (Einsteins formel $E = mc^2$), siden små deler av massen kan omdannes til energi.

\section{Forsøk A – Løselighet av salt i vann}
\subsection*{Utstyr}
\begin{itemize}
  \item Vekt (nøyaktighet $\pm 0.01$ g)
  \item Beger med lokk/folie
  \item Vann ($\approx 100$ g)
  \item Koksalt (NaCl, ca. 5 g)
\end{itemize}

\subsection*{Fremgangsmåte}
\begin{enumerate}
  \item Vei tomt beger med lokk: $m_\text{beger}$.
  \item Tilsett vann, vei på nytt: $m_\text{beger+vann}$.
  \item Vei saltet: $m_\text{salt}$.
  \item Løs opp saltet i vannet med lokk på.
  \item Vei begeret etter oppløsning: $m_\text{etter}$.
\end{enumerate}

\subsection*{Resultater}
\begin{table}[!ht]
\centering
\begin{tabular}{lcccccc}
\toprule
Prøve & $m_\text{beger+vann}$ (g) & $m_\text{salt}$ (g) & Forventet (g) & Målt (g) & Avvik (g) & Avvik (\%) \\
\midrule
A (salt i vann) & 150.00 & 5.00 & 155.00 & 154.99 & 0.01 & 0.0065 \\
\bottomrule
\end{tabular}
\end{table}

\subsection*{Beregning}
\[
\text{Avvik} = |m_\text{målt} - m_\text{forventet}|
\]
\[
\text{Prosentavvik} = \frac{\text{Avvik}}{m_\text{forventet}} \times 100\%
\]

\subsection*{Konklusjon}
Massetapet er neglisjerbart, og forsøket bekrefter loven om bevaring av masse.

\section{Forsøk B – Reaksjon mellom bakepulver og eddik}
\subsection*{Utstyr}
\begin{itemize}
  \item Reagensglass med ballong
  \item Vekt
  \item Eddik ($5-10mL$)
  \item Natron ($0,5-1g$)
\end{itemize}
\psset{unit=1cm}
\begin{pspicture}(-6,-2)(3,4)
\psset{unit=0.5}
  % Ballong (sirkel)
  \pscircle[linecolor=red!80!black,
            fillstyle=solid,
            fillcolor=red!40](0,3){2}

  % Utløp ballong
  \pspolygon[fillstyle=solid,fillcolor=red!40,
             linecolor=red!70!black]
             (-0.4,0.5)(-0.1,1.0)(0.4,1.0)(0.6,0.5)
\rput(0,-1.05){             
\pstTubeEssais[glassType=erlen,niveauLiquide1=80]
}
\psline[linecolor=red!40, linewidth=1.4pt](-0.1,1)(0.35 ,1)
  % Tekst (CO_2)
\rput(0,3){\large CO$_2$}
%Ring i ballong som fester rundt kolbe
\rput(0.056,0.42){  
\psellipse[linecolor=red!80!black,fillcolor=red!40,fillstyle=solid,linewidth=0.6pt](0.7,0.3)  
\psline[linecolor=red!40, linewidth=1.4pt](-0.4,0.1)(0.4 ,0.1)
}
\end{pspicture}
\subsection*{Fremgangsmåte}
\begin{enumerate}
  \item Vei flaske med eddik og ballong montert: $m_\text{før}$.
  \item Vei natron separat: $m_\text{natron}$.
  \item Hell natron ned i flasken ved å løfte ballongen.
  \item Etter reaksjonen: vei flaske + ballong igjen: $m_\text{etter}$.
\end{enumerate}

\subsection*{Resultater}
\begin{table}[h!]
\centering
\begin{tabular}{lcccccc}
\toprule
Prøve & $m_\text{flaske+eddik}$ (g) & $m_\text{natron}$ (g) & Forventet (g) & Målt (g) & Avvik (g) & Avvik (\%) \\
\midrule
B (natron + eddik) & 120.00 & 4.00 & 124.00 & 124.02 & 0.02 & 0.0161 \\
\bottomrule
\end{tabular}
\end{table}

\subsection*{Konklusjon}
Massen forblir tilnærmet konstant når gassen holdes inne i systemet. Forsøket bekrefter at massen bevares i kjemiske reaksjoner så lenge systemet er lukket.

\section{Feilkilder}
\begin{itemize}
  \item Lekkasje av gass i systemet
  \item Fordampning av vann
  \item Måleusikkerhet på vekten
  \item Tap av materiale under veiing eller overføring
\end{itemize}

\section{Videre arbeid}
\begin{itemize}
  \item Gjenta målingene flere ganger og ta gjennomsnitt.
  \item Undersøk hvordan temperatur påvirker nøyaktigheten.
  \item Diskuter kjernefysiske unntak der energi og masse henger sammen.
\end{itemize}


\section{Forsøk A – Salt løses i vann}

Dette er en \textbf{fysisk prosess}, ikke en kjemisk reaksjon. 
Når natriumklorid løses i vann, brytes de ioniske bindingene, og ionene stabiliseres av vannmolekylene:

\begin{align*}
\text{NaCl}_{(s)} &\rightarrow \text{Na}^{+}_{(aq)} + \text{Cl}^{-}_{(aq)}
\end{align*}

Her dannes ingen nye stoffer, og massen bevares fullstendig.

\section{Forsøk B – Natron og eddik}

Dette er en \textbf{kjemisk reaksjon} mellom natron (NaHCO$_3$) og eddiksyre (CH$_3$COOH). 
Reaksjonen danner natriumacetat, vann og karbondioksid:

\begin{align*}
\text{NaHCO}_{3(s)} + \text{CH}_{3}\text{COOH}_{(aq)} 
&\rightarrow \text{CH}_{3}\text{COONa}_{(aq)} + \text{H}_{2}\text{O}_{(l)} + \text{CO}_{2(g)}
\end{align*}

Reaksjonen skjer i to trinn:

\begin{align*}
\text{H}^{+} + \text{HCO}_{3}^{-} &\rightarrow \text{H}_{2}\text{CO}_{3} \\
\text{H}_{2}\text{CO}_{3} &\rightarrow \text{H}_{2}\text{O} + \text{CO}_{2}
\end{align*}

Dersom systemet er lukket, blir gassen værende, og total masse forblir konstant.
\maketitle  rapport

\section*{Instruks:}
Gå gjennom hvert punkt og kryss av når det er sjekket.

\section*{Forside / Tittel}
\begin{itemize}
    \item Tittel på forsøket tydelig
    \item Navn, dato, klasse og lærer
    \item Eventuelt fag og skole
\end{itemize}

\section*{Hensikt / Problemstilling}
\begin{itemize}
    \item Hensikten med forsøket tydelig formulert
    \item Koblet til loven om bevaring av masse
\end{itemize}

\section*{Teori}
\begin{itemize}
    \item Kort forklaring av loven om bevaring av masse
    \item Forskjell på fysisk og kjemisk prosess
    \item Historisk bakgrunn (valgfritt)
\end{itemize}

\section*{Utstyr}
\begin{itemize}
    \item Alle materialer listet
    \item Merk om lukkede systemer brukes (lokk/ballong)
\end{itemize}

\section*{Fremgangsmåte}
\begin{itemize}
    \item Trinnvis og tydelig
    \item Lett for andre å gjenta forsøket
\end{itemize}

\section*{Kjemiske reaksjoner}
\begin{itemize}
    \item Reaksjonslikning for Forsøk A (salt i vann)
    \item Reaksjonslikning for Forsøk B (natron + eddik)
    \item Tilstandsangivelser inkludert (s, l, aq, g)
    \item Kort forklaring av hva som skjer
\end{itemize}

\section*{Resultater}
\begin{itemize}
    \item Tabell med målinger: forventet masse, målt masse
    \item Absolutt og prosentavvik
    \item Kommentar om måleusikkerhet
\end{itemize}

\section*{Beregninger}
\begin{itemize}
    \item Avvik og prosentavvik beregnet og vist
    \item Tall stemmer med tabell
\end{itemize}

\section*{Diskusjon / Konklusjon}
\begin{itemize}
    \item Konkluder med at masse er bevart
    \item Drøft små avvik (fordampning, lekkasje, vektusikkerhet)
    \item Refleksjon: betydning av lokk/ballong
\end{itemize}

\section*{Feilkilder / Forbedringer}
\begin{itemize}
    \item Kort liste med realistiske feilkilder
    \item Forslag til forbedringer
\end{itemize}

\section*{Språk / Formatering her}
\begin{itemize}
    \item Rettstavning og tegnsetting sjekket
    \item Konsistent bruk av symboler og enheter (g, ml)
\end{itemize}
\begin{tabular}{l l l p{3cm}}
\toprule
\textbf{Beholder} & \textbf{Eddiksyre (mL)} & \textbf{Natron (g)} & \textbf{Kommentar} \\
\midrule
Reagensglass & 5--10 & 0,5--1 & Små mengder, lett å feste ballong, liten plass til gass \\
Glassflaske / glassbeholder & 20--50 & 2--5 & Større volum, oversiktlig reaksjon, ballong eller lokk kan brukes \\
Plastflaske & 20--50 & 2--5 & Praktisk for ballong og vektmåling, fleksibel og trygg \\
\bottomrule
\end{tabular}
\end{document}